\documentclass[11pt,a4paper]{scrartcl}

\usepackage{amssymb}
\usepackage{tikz} 
\usepackage{amsmath}

\usepackage[utf8]{inputenc}
\usepackage[T1]{fontenc}
\newcommand{\ats}[1]{\par\noindent\textit{Atsakymas:} #1}
\usepackage{cancel}
\usepackage{tikz}
\usepackage{lmodern} 
\usepackage{amsmath,amssymb,amsthm}
\usepackage{graphicx}
\usepackage[dvipsnames,svgnames]{xcolor}
\usepackage{hyperref}
\usepackage{mathrsfs}
\usepackage[shortlabels]{enumitem}
\usepackage[obeyFinal,textsize=scriptsize,shadow,loadshadowlibrary]{todonotes}
\usepackage{mathtools}
\usepackage{microtype}

%%% Paragraph layout
\setlength{\parindent}{0pt}
\setlength{\parskip}{0.6\baselineskip}

%%% Colored sections
\renewcommand*{\sectionformat}%
  {\color{purple}\S\thesection\enskip}
\renewcommand*{\subsectionformat}%
  {\color{purple}\S\thesubsection\enskip}
\renewcommand*{\subsubsectionformat}%
  {\color{purple}\S\thesubsubsection\enskip}
\addtokomafont{paragraph}{\color{orange!35!black}\P\ }
\KOMAoptions{numbers=noenddot}

%%% Title
\addtokomafont{subtitle}{\Large}
\setkomafont{author}{\Large\scshape}
\setkomafont{date}{\Large\normalsize}
% Neigiamas vertikalus tarpas, kuris pakelia dokumento antraštę į viršų. 
% Galite keisti -2.5cm į -3cm (aukščiau) arba -1.5cm (žemiau).
\titlehead{\vspace*{-7.5cm}} 

% PRIDĖTA: Automatiškai perrašome datą į mokyklos pavadinimą
\AtBeginDocument{%
  \date{\normalsize Vilniaus licėjus}
}

%%% Page setup
\usepackage[headsepline]{scrlayer-scrpage}
\renewcommand{\headfont}{}
\addtolength{\textheight}{3.14cm}
\setlength{\footskip}{0.5in}
\setlength{\headsep}{10pt}
% Puslapio antraštė turi rodyti tik konkrečius dokumento duomenis.
% Nustatome juos aiškiai, kad neatsirastų "\@author" / "\@title" arba senų reikšmių.
\makeatletter
\AtBeginDocument{%
  \ihead{\footnotesize\textbf{Andrius Berniukevičius}}%
  \ohead{\footnotesize\textbf{\@title}}%
}
\makeatother
\automark{section}
\chead{}
\cfoot{\pagemark}

%%% Macros
\providecommand{\ol}{\overline}
\providecommand{\ul}{\underline}
\providecommand{\wt}{\widetilde}
\providecommand{\wh}{\widehat}
\providecommand{\eps}{\varepsilon}
\providecommand{\half}{\frac{1}{2}}
\providecommand{\inv}{^{-1}}
\newcommand{\dang}{\measuredangle} %% Directed angle
\newcommand{\vocab}[1]{\textbf{\color{blue}\sffamily #1}}
\providecommand{\CC}{\mathbb C}
\providecommand{\FF}{\mathbb F}
\providecommand{\NN}{\mathbb N}
\providecommand{\QQ}{\mathbb Q}
\providecommand{\RR}{\mathbb R}
\providecommand{\ZZ}{\mathbb Z}
\providecommand{\ts}{\textsuperscript}
\providecommand{\dg}{^\circ}
\providecommand{\ii}{\item}
\providecommand{\defeq}{\coloneqq}
\DeclareMathOperator*{\lcm}{lcm}
\DeclareMathOperator*{\argmin}{arg min}
\DeclareMathOperator*{\argmax}{arg max}
\providecommand{\hrulebar}{\par
\hspace{\fill}\rule{0.95\linewidth}{.7pt}\hspace{\fill}
\par\nointerlineskip \vspace{\baselineskip}}

%%% Optional colored theorems
\usepackage{thmtools}
\usepackage[framemethod=TikZ]{mdframed}
\mdfdefinestyle{mdbluebox}{%
  roundcorner=10pt,
  linewidth=1pt,
  skipabove=12pt,
  innerbottommargin=9pt,
  skipbelow=2pt,
  linecolor=blue,
  nobreak=true,
  backgroundcolor=TealBlue!5,
}
\declaretheoremstyle[
  headfont=\sffamily\bfseries\color{MidnightBlue},
  mdframed={style=mdbluebox},
  headpunct={\\[3pt]},
  postheadspace={0pt}
]{thmbluebox}
\mdfdefinestyle{mdredbox}{%
  linewidth=0.5pt,
  skipabove=12pt,
  frametitleaboveskip=5pt,
  frametitlebelowskip=0pt,
  skipbelow=2pt,
  frametitlefont=\bfseries,
  innertopmargin=4pt,
  innerbottommargin=8pt,
  nobreak=true,
  backgroundcolor=Salmon!5,
  linecolor=RawSienna,
}
\declaretheoremstyle[
  headfont=\bfseries\color{RawSienna},
  mdframed={style=mdredbox},
  headpunct={\\[3pt]},
  postheadspace={0pt},
]{thmredbox}
\mdfdefinestyle{mdgreenbox}{%
  skipabove=8pt,
  linewidth=2pt,
  rightline=false,
  leftline=true,
  topline=false,
  bottomline=false,
  linecolor=ForestGreen,
  backgroundcolor=ForestGreen!5,
}
\declaretheoremstyle[
  headfont=\bfseries\sffamily\color{ForestGreen!70!black},
  bodyfont=\normalfont,
  spaceabove=2pt,
  spacebelow=1pt,
  mdframed={style=mdgreenbox},
  headpunct={ --- },
]{thmgreenbox}
\mdfdefinestyle{mdblackbox}{%
  skipabove=8pt,
  linewidth=3pt,
  rightline=false,
  leftline=true,
  topline=false,
  bottomline=false,
  linecolor=black,
  backgroundcolor=RedViolet!5!gray!5,
}
\declaretheoremstyle[
  headfont=\bfseries,
  bodyfont=\normalfont\small,
  spaceabove=0pt,
  spacebelow=0pt,
  mdframed={style=mdblackbox}
]{thmblackbox}

%% Aplinkos su rėmeliais (thmtools)
\declaretheorem[style=thmbluebox,name=Teorema]{theorem}
\declaretheorem[style=thmbluebox,name=Lema]{lemma}
\declaretheorem[style=thmbluebox,name=Savybė]{property}
\declaretheorem[style=thmbluebox,name=Teiginys]{proposition}
\declaretheorem[style=thmbluebox,name=Išvada]{corollary}
\declaretheorem[style=thmbluebox,name=Teorema,numbered=no]{theorem*}
\declaretheorem[style=thmbluebox,name=Lema,numbered=no]{lemma*}
\declaretheorem[style=thmbluebox,name=Teiginys,numbered=no]{proposition*}
\declaretheorem[style=thmbluebox,name=Išvada,numbered=no]{corollary*}
\declaretheorem[style=thmgreenbox,name=Algoritmas]{algorithm}
\declaretheorem[style=thmgreenbox,name=Algoritmas,numbered=no]{algorithm*}
\declaretheorem[style=thmgreenbox,name=Teiginys]{claim}
\declaretheorem[style=thmgreenbox,name=Teiginys,numbered=no]{claim*}
\declaretheorem[style=thmredbox,name=Pavyzdys]{example}
\declaretheorem[style=thmredbox,name=Pavyzdys,numbered=no]{example*}
\declaretheorem[style=thmblackbox,name=Pastaba]{remark}
\declaretheorem[style=thmblackbox,name=Pastaba,numbered=no]{remark*}

%% Standartinės amsthm aplinkos (Apibrėžimai ir užduotys)
\theoremstyle{definition}
\ifcsname conjecture\endcsname\else\newtheorem{conjecture}{Hipotezė}\fi
\ifcsname definition\endcsname\else\newtheorem{definition}{Apibrėžimas}\fi
\ifcsname fact\endcsname\else\newtheorem{fact}{Faktas}\fi
\ifcsname answer\endcsname\else\newtheorem{answer}{Atsakymas}\fi
\ifcsname ques\endcsname\else\newtheorem{ques}{Klausimas}\fi
\ifcsname exercise\endcsname\else\newtheorem{exercise}{Užduotis}\fi
\ifcsname problem\endcsname\else\newtheorem{problem}{Uždavinys}[section]\fi
\ifcsname conjecture*\endcsname\else\newtheorem*{conjecture*}{Hipotezė}\fi
\ifcsname definition*\endcsname\else\newtheorem*{definition*}{Apibrėžimas}\fi
\ifcsname fact*\endcsname\else\newtheorem*{fact*}{Faktas}\fi
\ifcsname answer*\endcsname\else\newtheorem*{answer*}{Atsakymas}\fi
\ifcsname ques*\endcsname\else\newtheorem*{ques*}{Klausimas}\fi
\ifcsname exercise*\endcsname\else\newtheorem*{exercise*}{Užduotis}\fi
\ifcsname problem*\endcsname\else\newtheorem*{problem*}{Uždavinys}\fi

\newcounter{answerssection}
\newcommand{\answerssection}{%
  \setcounter{answerssection}{\value{section}}%
  \section{Atsakymai}%
  \setcounter{problem}{0}%
  \renewcommand{\theproblem}{\arabic{answerssection}.\arabic{problem}}%
}

%% GALUTINIS SPRENDIMAS PROOF APLINKAI
%% --- PRADŽIA: Įrodymo aplinkos sutvarkymas ---
\makeatletter

% 1. Užtikriname, kad žodis būtų "Įrodymas", net jei "babel" paketas bando jį perrašyti
\AtBeginDocument{%
  \renewcommand{\proofname}{Įrodymas}%
  \@ifpackageloaded{babel}{%
    \addto\captionslithuanian{\renewcommand{\proofname}{Įrodymas}}%
  }{}%
}

% 2. Priverstinai pašaliname scrartcl klasės bold šriftą ir paliekame TIK italic
% 3. Sujungiame su tuščiu kvadratėliu dešiniajame kampe.
\renewcommand{\qedsymbol}{\ensuremath{\Box}}
\renewenvironment{proof}[1][\proofname]{\par
  \pushQED{\qed}%
  \normalfont \topsep6\p@\@plus6\p@\relax
  \trivlist
  \item[\hskip\labelsep
        \normalfont\itshape % Priverčia naudoti tik pasvirą šriftą, kaip nuotraukoje
    #1\@addpunct{.}]\ignorespaces
}{%
  \popQED\endtrivlist\@endpefalse
}
\newenvironment{solution}{\par
  \normalfont \topsep6\p@\@plus6\p@\relax
  \trivlist
  \item[\hskip\labelsep
        \normalfont\itshape
    \textit{Sprendimas}\@addpunct{.}]\ignorespaces
}{%
  \endtrivlist\@endpefalse
}
\newcommand{\ans}[1]{\par\noindent\textit{Atsakymas:} #1}
\makeatother


\title{Dirichlė (Balandžių ir narvelių) principas}
\author{Andrius Berniukevičius}

\newenvironment{solution}
    {\begin{list}{}{\setlength{\leftmargin}{10pt}\setlength{\rightmargin}{10pt}}\item[]}
    {\end{list}}

\begin{document}

\maketitle

\section{Kas yra Dirichlė principas?}

Dirichlė (dažnai vadinamas \textit{balandžių ir narvelių}) principas yra vienas intuityviausių, tačiau galingiausių įrankių sprendžiant logikos, kombinatorikos ir skaičių teorijos uždavinius. Skirtingai nei invariantai, rodantys, kad procesas neįmanomas, Dirichlė principas įrodo, kad tam tikra situacija \textbf{garantuotai egzistuoja}.

\begin{theorem}
Jei $(n+1)$ balandžių yra patalpinami į $n$ narvelių, tai bent viename narvelyje bus \textbf{bent du balandžiai}.
\end{theorem}

\begin{theorem}[Apibendrintas]
Tegu $m$ ir $n$ yra teigiami sveikieji skaičiai. Jei $(mn + 1)$ balandžių patalpinami į $n$ narvelių, tai bent viename narvelyje garantuotai atsidurs \textbf{bent $(m + 1)$ balandžių}.
\end{theorem}

Sprendžiant uždavinius sunkiausia yra kūrybiškai sugalvoti, kas konkrečioje situacijoje atlieka \textbf{„balandžių“}, o kas – \textbf{„narvelių“} vaidmenį. „Narveliai“ dažniausiai būna dalybos liekanos, skaičių poros, kurių suma fiksuta, arba tam tikros geometrinės figūros dalys.

\subsection{Pavyzdžiai}

\begin{example}
Įrodykite, kad iš bet kokių keturių sveikųjų skaičių visada galima išrinkti du, kurių skirtumas dalijasi iš 3.
\end{example}

\begin{solution}
\textbf{Sprendimas:} Bet kokį skaičių dalijant iš 3, galimos tik trys liekanos: 0, 1 arba 2. Taigi, mes turime 3 „narvelius“ (liekanas) ir 4 „balandžius“ (skaičius). Pagal Dirichlė principą, bent du skaičiai turės tą pačią liekaną. Kadangi jų liekanos vienodos, šių skaičių skirtumas dalysis iš 3. \hfill $\square$
\end{solution}

\begin{example}
Kvadrate, kurio kraštinės ilgis lygus 2, atsitiktinai pažymėti 5 taškai. Įrodykite, kad atsiras bent du taškai, atstumas tarp kurių yra ne didesnis nei $\sqrt{2}$.
\end{example}

\begin{solution}
\textbf{Sprendimas:} Padalykime didįjį kvadratą į 4 mažesnius kvadratus (narvelius), kurių kiekvieno kraštinė yra 1.

\begin{center}
\begin{tikzpicture}[scale=1.5]
    \draw[thick] (0,0) rectangle (2,2);
    \draw[thick] (1,0) -- (1,2);
    \draw[thick] (0,1) -- (2,1);
    \node at (0.5, 0.5) {$1 \times 1$};
    \node at (1.5, 0.5) {$1 \times 1$};
    \node at (0.5, 1.5) {$1 \times 1$};
    \node at (1.5, 1.5) {$1 \times 1$};
    \filldraw[black] (0.3,0.7) circle (1.5pt);
    \filldraw[black] (1.7,0.2) circle (1.5pt);
    \filldraw[black] (1.2,1.8) circle (1.5pt);
    \filldraw[black] (0.8,1.3) circle (1.5pt);
    \filldraw[black] (0.6,1.8) circle (1.5pt);
\end{tikzpicture}
\end{center}

Kadangi turime 5 taškus (balandžius) ir 4 kvadratus (narvelius), bent du taškai atsidurs tame pačiame $1 \times 1$ kvadrate. Didžiausias įmanomas atstumas tarp dviejų taškų tokiame kvadrate yra jo įstrižainė, kurios ilgis pagal Pitagoro teoremą lygus $\sqrt{1^2 + 1^2} = \sqrt{2}$. \hfill $\square$
\end{solution}

\begin{example}
Iš skaičių aibės $\{1, 2, 3, \dots, 100\}$ atsitiktinai išrenkamas 51 skaičius. Įrodykite, kad tarp jų būtinai atsiras du skaičiai, iš kurių vienas dalijasi iš kito.
\end{example}

\begin{solution}
\textbf{Sprendimas:} Bet kokį sveikąjį skaičių unikaliai galime užrašyti pavidalu $j \cdot 2^k$, kur $j$ – nelyginis skaičius, o $k \ge 0$. Aibėje nuo 1 iki 100 yra lygiai 50 nelyginių skaičių ($1, 3, 5, \dots, 99$). Šias nelygines dalis naudosime kaip „narvelius“. 
Kadangi išrenkame 51 skaičių, bent du iš jų turės tą pačią nelyginę dalį $j$. Taigi, vienas iš tų skaičių bus $j \cdot 2^{k_1}$, o kitas $j \cdot 2^{k_2}$. Akivaizdu, kad tas skaičius, kurio dvejeto laipsnis yra mažesnis, dalys tą skaičių, kurio dvejeto laipsnis didesnis. \hfill $\square$
\end{solution}


\newpage
\section{Uždaviniai savarankiškam sprendimui (30 iššūkių)}

\begin{enumerate}
    \renewcommand{\labelenumi}{\textbf{\theenumi.}}
    
    % --- Lengvi (Klasikiniai balandžiai) ---
    \item Mokykloje mokosi 367 mokiniai. Įrodykite, kad bent du iš jų švenčia gimtadienį tą pačią dieną.
    \item Miške auga vienas milijonas pušų. Yra žinoma, kad ant jokios pušies nėra daugiau nei $600\,000$ spyglių. Įrodykite, kad miške garantuotai yra bent dvi pušys, turinčios lygiai tiek pat spyglių.
    \item Tamsiame kambaryje stalčiuje guli 10 juodų ir 10 baltų kojinių. Kiek mažiausiai kojinių reikia ištraukti nežiūrint, kad garantuotai turėtumėte bent vieną vienodos spalvos porą?
    \item Dėžėje yra daug raudonų, mėlynų, žalių ir geltonų rutuliukų. Kiek mažiausiai rutuliukų reikia ištraukti aklai, kad garantuotai turėtumėte bent \textbf{tris} tos pačios spalvos rutuliukus?
    \item Iš aibės $\{1, 2, 3, \dots, 10\}$ atsitiktinai išrenkami 6 skaičiai. Įrodykite, kad tarp jų būtinai atsiras du, kurių suma lygi 11.
    \item Iš standartinės 52 kortų malkos aklai traukiamos kortos. Kiek mažiausiai kortų reikia ištraukti, kad rankose garantuotai turėtumėte bent 5 tos pačios rūšies (širdžių, būgnų, pikų arba kryžių) kortas?
    
    % --- Vidutiniai (Kombinatorika ir liekanos) ---
    \item Gimtadienio vakarėlyje dalyvauja $n$ žmonių. Įrodykite, kad vakarėlyje visada atsiras bent du žmonės, turintys lygiai tokį patį skaičių pažįstamų tarp kitų vakarėlio svečių.
    \item Matematikos teste yra tik 2 uždaviniai. Už teisingai išspręstą skiriama 10 taškų, už praleistą – 2 taškai, už neteisingą – 0 taškų. Kiek mažiausiai mokinių turi laikyti testą, kad bent 3 iš jų surinktų vienodą galutinę taškų sumą?
    \item Dėžėje guli daug keturių skirtingų rūšių Velykinių kiaušinių. Kiekvienam vaikui leidžiama iš dėžės išsitraukti po 1 arba 2 kiaušinius. Kiek mažiausiai vaikų turi dalyvauti žaidime, kad bent du vaikai padarytų visiškai identišką pasirinkimą?
    \item Apskritime surašyti 10 sveikųjų skaičių, kurių suma lygi 100. Įrodykite, kad visada atsiras trys greta esantys skaičiai, kurių suma yra ne mažesnė kaip 30.
    \item Turime bet kokius $(n+1)$ sveikųjų skaičių. Įrodykite, kad visada galima išrinkti du skaičius, kurių skirtumas dalijasi iš $n$.
    \item Duota 12 skirtingų dviženklių skaičių. Įrodykite, kad iš jų visada galima išrinkti du skaičius taip, kad jų skirtumo vienetų ir dešimčių skaitmenys būtų vienodi (t. y. skirtumas dalytųsi iš 11).
    \item Duoti bet kokie 5 sveikieji skaičiai. Įrodykite, kad iš jų visada galima išrinkti lygiai 3 skaičius taip, kad jų suma dalytųsi iš 3.
    \item Plokštumoje pažymėti 5 taškai, kurių visos koordinatės $(x, y)$ yra sveikieji skaičiai. Įrodykite, kad bent vienos atkarpos, jungiančios bet kuriuos du iš šių taškų, vidurio taškas taip pat turi tik sveikas koordinates.
    \item Duota 100 bet kokių sveikųjų skaičių. Įrodykite, kad iš jų galima išrinkti bent 15 skaičių taip, kad bet kurių dviejų išrinktųjų skirtumas dalytųsi iš 7.
    \item Tegu $a_1, a_2, \dots, a_9$ yra skaičių $1, 2, \dots, 9$ kėlinys (tie patys skaičiai surašyti bet kokia tvarka). Įrodykite, kad sandauga $(a_1-1)(a_2-2)\dots(a_9-9)$ būtinai yra lyginis skaičius.
    
    % --- Geometriniai iššūkiai ---
    \item Kvadrate, kurio matmenys yra $3 \times 3$, atsitiktinai pažymėta 10 taškų. Įrodykite, kad atsiras bent du taškai, atstumas tarp kurių yra ne didesnis nei $\sqrt{2}$.
\begin{center}
\begin{tikzpicture}[scale=0.8]
    \draw[thick] (0,0) rectangle (3,3);
    \draw[step=1cm,gray,very thin,dashed] (0,0) grid (3,3);
\end{tikzpicture}
\end{center}
    \item Lygiakraščio trikampio, kurio kraštinės ilgis yra 2, viduje atsitiktinai pažymėti 5 taškai. Įrodykite, kad atstumas tarp kažkurių dviejų taškų neviršija 1.
\begin{center}
\begin{tikzpicture}[scale=1.2]
    \draw[thick] (0,0) -- (2,0) -- (1, {sqrt(3)}) -- cycle;
    \draw[dashed, gray] (1,0) -- (0.5, {sqrt(3)/2}) -- (1.5, {sqrt(3)/2}) -- cycle;
\end{tikzpicture}
\end{center}
    \item Taisyklingojo šešiakampio, kurio kraštinės ilgis lygus 1, viduje atsitiktinai pažymėti 7 taškai. Įrodykite, kad atsiras bent du taškai, atstumas tarp kurių neviršija 1.
    \item Skritulio, kurio spindulys yra 1, viduje atsitiktinai pažymėti 7 taškai. Įrodykite, kad atsiras bent du taškai, atstumas tarp kurių neviršija 1.
    \item Vienetinio kvadrato ($1 \times 1$) viduje pažymėti 9 taškai. Įrodykite, kad iš jų visada galima išrinkti 3 taškus, kurie sudarytų trikampį, kurio plotas neviršytų $1/8$.
    
    % --- Olimpiadiniai „Bossai“ ---
    \item Duoti bet kokie 52 sveikieji skaičiai. Įrodykite, kad iš jų visada galima išrinkti du skaičius taip, kad tų skaičių suma arba skirtumas dalytųsi iš 100.
    \item Duota 7 atsitiktinių sveikųjų skaičių seka. Įrodykite, kad visada įmanoma rasti vieną arba kelis \textit{iš eilės einančius} šios sekos elementus, kurių suma dalijasi iš 7.
    \item Įrodykite, kad egzistuoja skaičiaus 2017 kartotinis, kurio dešimtainiame užraše naudojami išskirtinai tik skaitmenys 8 ir 0.
    \item Fibonači sekoje $1, 1, 2, 3, 5, 8, \dots$ kiekvienas narys lygus dviejų ankstesnių narių sumai. Įrodykite, kad šioje sekoje privalo egzistuoti skaičius, kuris baigiasi bent keturiais nuliais.
    \item Visi begalinės plokštumos taškai yra atsitiktinai nuspalvinti dviem spalvomis (pavyzdžiui, raudona ir mėlyna). Įrodykite, kad plokštumoje būtinai atsiras du \textbf{vienodos spalvos} taškai, atstumas tarp kurių yra lygiai 1.
    \item Kiekvienas plokštumos taškas su sveikomis koordinatėmis $(x, y)$ yra nuspalvintas viena iš 3 spalvų. Įrodykite, kad visada egzistuoja stačiakampis, kurio kraštinės lygiagrečios koordinačių ašims, o visos keturios viršūnės yra tos pačios spalvos.
    \item Yra dvi skirtingo dydžio ruletės su bendru centru. Kiekviena jų padalinta į 200 vienodų sektorių. Didesnėje ruletėje 100 sektorių nuspalvinti raudonai, kiti 100 – žaliai. Mažesnės ruletės 200 sektorių atsitiktinai nuspalvinti raudonai arba žaliai (spalvų skaičius nebūtinai lygus). Įrodykite, kad visada galima pasukti ruletes taip, jog bent 100 persidengiančių sektorių poros sutaptų spalvomis.
    \item Lentoje surašytas 101 atsitiktinis skirtingas skaičius. Įrodykite, kad iš jų visada galima išbraukti 90 skaičių taip, kad likę 11 skaičių eitų griežtai didėjimo arba griežtai mažėjimo tvarka.
    \item \textbf{(Kombinatorikos finalas)} Begalinėje languotoje plokštumoje kiekvienas langelis nuspalvintas viena iš $N$ spalvų. Įrodykite, kad visada galima rasti keturis tos pačios spalvos langelius, kurių centrai sudaro stačiakampį su kraštinėmis, lygiagrečiomis tinklelio linijoms.
\end{enumerate}

\end{document}